% Part: sets-functions-relations
% Chapter: size-of-sets-complete

\documentclass[../../../include/open-logic-chapter]{subfiles}

\begin{document}

\olchapter{sfr}{siz}{سیٹاں دا سائز (The Size of Sets)}

\begin{editorial}
ایہ باب گنتی وار فہرستاں (enumerations)، قابلِ شمار ہون
(countability) تے ناقابلِ شمار ہون (uncountability) اُتے گل کردا اے۔
کئی حصے دو روپاں وچ دتے گئے نیں: اک ہور بنیادی روپ (more elementary
version)، جیہڑا گنتی وار فہرستاں نوں لسٹاں (lists)، یا $\PosInt$ توں
ہر ہدف قدر لین والے فنکشناں (surjections) دے طور تے لیندا اے؛ تے اک
ہور مجرد روپ (more abstract version)، جیہڑا گنتی وار فہرستاں نوں
$\Nat$ نال اک بہ اک مطابقتاں (bijections) دے طور تے تعریف کردا اے۔
\end{editorial}

\olimport{introduction}

\olimport{enumerability}

\olimport{zig-zag}

\olimport{pairing}

\olimport{pairing-alt}

\olimport{non-enumerability}

\olimport{reduction}

\olimport{equinumerous-sets}

\olimport{comparing-size}

\olimport{schroder-bernstein}

\begin{editorial}
اگلے \olref[sfr][siz][enm-alt]{sec}،
\olref[sfr][siz][nen-alt]{sec} تے \olref[sfr][siz][red-alt]{sec}،
\olref[sfr][siz][enm]{sec}، \olref[sfr][siz][nen]{sec} تے
\olref[sfr][siz][red]{sec} دے متبادل روپ (alternative versions) نیں۔
ایہ ٹم بٹن نے اپنی
اوپن سیٹ تھیوری (Open Set Theory) دی کتاب وچ ورتن لئی لکھے سن۔ ایہ
تھوڑے ہور اگے دے درجے دے (slightly more advanced) نیں تے گنتی وار
فہرست بنن جوگ ہون (enumerability) دی اک وکھری تعریف ورتدے نیں، جیہڑی
سیٹ تھیوری دے سیاق (set-theory context) وچ ہور مناسب اے: یعنی $\Nat$
نال یا اوہدے کسے مُڈھلے ٹکڑے (initial segment) نال اک بہ اک مطابقت
(bijection)، نہ کہ فہرست بن سکنا (being listable) یا $\PosInt$ توں ہر
ہدف قدر لین والے فنکشن (surjective function) دا حاصل قدراں دا سیٹ
(range) ہونا۔
% PNBSIZ-001 مقامی تدوینی درستی: ماخذ دے "a difference definition" نوں سیاق مطابق "a different definition" (اک وکھری تعریف) کیتا؛ منجمد ماخذ نہیں بدلیا۔
\end{editorial}

\olimport{enumerability-alt}
\olimport{non-enumerability-alt}
\olimport{reduction-alt}

\OLEndChapterHook

\end{document}
